\section{Support geometry of the covariance kernel}\label{sec:support84}
Let $\mathcal H_d$ be the real Hilbert space of Hermitian $d\times d$
matrices, with inner product $\tr(AB)$.  Throughout this section,
$E=(E_1,\ldots,E_k)$ is an arbitrary ordered measurement:
$E_j\succeq0$ and $\sum_jE_j=I_d$.  Put
\[
 P_j=\operatorname{supp}E_j,\qquad Q_j=I_d-P_j,\qquad r_j=\operatorname{rank}E_j.
\]
The support of a zero effect is the zero projection.  The real subspaces
\begin{equation}\label{eq:supportspan84}
 \mathcal W_E=\sum_{j=1}^k P_j\mathcal H_dP_j,
 \qquad \mathcal A_E=\mathcal W_E^\perp
   =\{A\in\mathcal H_d:P_jAP_j=0\text{ for every }j\}
\end{equation}
will be called the \emph{support span} and its orthogonal complement.
The sum in \eqref{eq:supportspan84} is a sum of real linear subspaces,
not an assertion that $\mathcal W_E$ is an algebra.  In particular,
for higher-rank effects it need not equal the span of the effects themselves.
We use the zero-sum tuple space
$\mathcal T=\{K\in\mathcal H_d^k:\sum_jK_j=0\}$ with the sum
Hilbert--Schmidt inner product, and the covariance of
Lemma~\ref{lem:covarianceidentities83}.

\begin{theorem}[Complete support description of the kernel]\label{thm:supportkernel84}
For $A\in\mathcal A_E$ and $Z_j\in Q_j\mathcal H_dQ_j$, set
\begin{equation}\label{eq:kernelparam84}
 X_j=i[P_j,A]+Z_j,\qquad
 K_j=X_j-\frac1k\sum_{\ell=1}^kX_\ell.
\end{equation}
This is a real linear bijection from
$\mathcal A_E\oplus\bigoplus_jQ_j\mathcal H_dQ_j$ onto
$\ker(\mathcal C_E|_{\mathcal T})$.  Consequently
\begin{equation}\label{eq:kerneldimension84}
 \dim\ker(\mathcal C_E|_{\mathcal T})
 =d^2-\dim\mathcal W_E+\sum_j(d-r_j)^2.
\end{equation}
In particular the kernel, and not just its dimension, depends only
on the ordered support projections.
\end{theorem}
\begin{proof}
Write $S=\sum_jE_jK_j$.  The nonnegative variance decomposition in
Lemma~\ref{lem:covarianceidentities83} shows that $K$ is in the kernel
exactly when
\begin{equation}\label{eq:supportkernelcondition84}
 P_jK_j=P_jS\quad\text{for every }j.
\end{equation}
Indeed multiplication by $\sqrt{E_j}$ has precisely the same kernel
as multiplication by $P_j$; no inverse at a singular effect is used.
Decompose $S=B+iA$, where $A,B$ are Hermitian.  Since
$P_jK_jP_j$ and $P_jBP_j$ are Hermitian, compressing
\eqref{eq:supportkernelcondition84} gives $P_jAP_j=0$.
Furthermore $P_j(K_j-B)=iP_jA$, with the adjoint identity on the
other side.  These determine the two off-diagonal support blocks
and the support block of $K_j-B$; its remaining block is arbitrary.
Thus
\[
 K_j=B+i[P_j,A]+Z_j,\qquad Z_j=Q_jZ_jQ_j=Z_j^*.
\]
The condition $\sum_jK_j=0$ forces
$B=-k^{-1}\sum_j(i[P_j,A]+Z_j)$, proving necessity.

Conversely, let $A,Z_j$ be as stated.  Because $P_jAP_j=0$,
$E_jAP_j=0$, and hence
\[
 \sum_jE_jX_j=i\sum_jE_jA=iA.
\]
For the mean-subtracted tuple, therefore, $S=iA-\overline X$,
where $\overline X=k^{-1}\sum_jX_j$ is Hermitian.
We have $P_jX_j=iP_jA$, which proves
\eqref{eq:supportkernelcondition84}.  To check injectivity, suppose
that all $K_j$ vanish.  Then all $X_j$ equal a common Hermitian
matrix $D$, so the last display gives $D=iA$.  A matrix that is
both Hermitian and anti-Hermitian is zero.  Therefore $A=D=0$ and
all $Z_j=0$.  Finally $\dim(Q_j\mathcal H_dQ_j)=(d-r_j)^2$;
rank--nullity in \eqref{eq:supportspan84} proves
\eqref{eq:kerneldimension84}.
\end{proof}

\begin{proposition}[Unitary directions and the support span]\label{prop:unitaryrange84}
Let $G\in\mathcal H_d$ and $H_j=i[G,E_j]$.  Then
\begin{equation}\label{eq:unitaryrange84}
 H\in\operatorname{Ran}(\mathcal C_E|_{\mathcal T})
 \quad\Longleftrightarrow\quad G\in\mathcal W_E.
\end{equation}
More precisely, the pairing of $H$ with the kernel tuple
\eqref{eq:kernelparam84} is $-2\tr(GA)$.
\end{proposition}
\begin{proof}
We have $Q_jH_jQ_j=0$, so the $Z_j$ terms have zero pairing with $H$;
its zero tuple sum also removes $\overline X$.  Trace cyclicity gives
\[
 \sum_j\tr\bigl(H_ji[P_j,A]\bigr)
 =\tr\left(A\sum_ji[H_j,P_j]\right).
\]
Expanding in the displayed order yields
\[
 \sum_ji[H_j,P_j]
 =-2G+\sum_j(E_jGP_j+P_jGE_j).
\]
Every summand in the last sum belongs to
$P_j\mathcal H_dP_j$; it is therefore orthogonal to $A$.
The pairing is $-2\tr(GA)$ as asserted.  A self-adjoint linear
operator in finite dimension has range equal to the orthogonal
complement of its kernel.  Theorem~\ref{thm:supportkernel84} now
proves \eqref{eq:unitaryrange84}.
\end{proof}

\begin{corollary}[Regularized tangent scales]\label{cor:tangentscales84}
Let $\Pi_0$ be the orthogonal projection onto
$\ker(\mathcal C_E|_{\mathcal T})$, and let $u_\ell$ be an orthonormal
basis of its complement with $\mathcal C_Eu_\ell=\lambda_\ell u_\ell$,
$\lambda_\ell>0$.  For every $H\in\mathcal T$,
\begin{equation}\label{eq:tangentspectral84}
 g_{N,E}(H)^2
 =N^2\|\Pi_0H\|_{\mathrm{HS}}^2+
  N\sum_\ell\frac{|\langle H,u_\ell\rangle|^2}{\lambda_\ell+N^{-1}}.
\end{equation}
Thus a nonzero horizontal direction has order $\sqrt N$, whereas a
direction with nonzero kernel component has order $N$, at a fixed
measurement.  For $H=i[G,E]$ the first alternative is exactly
$G\in\mathcal W_E$ and $H\ne0$.
\end{corollary}
\begin{proof}
Apply the finite-dimensional spectral theorem to
$N(\mathcal C_E+N^{-1}\operatorname{Id})^{-1}$ and use
Proposition~\ref{prop:unitaryrange84}.
\end{proof}
The corollary concerns a comparison form, not yet an operational
lower bound.  The distinction matters at the boundary; the next
section proves the corresponding finite-pair statement on entire
unitary orbits.

\begin{proposition}[Exact represented-input decision]\label{prop:supportexact84}
On Gaussian-rational effects and a Gaussian-rational Hermitian
generator, membership in $\mathcal W_E$, stationarity of the orbit,
and the kernel dimension \eqref{eq:kerneldimension84} are decidable
by exact rational linear algebra in polynomial bit complexity.
The same calculation gives the orthogonal projection of $G$ onto
$\mathcal W_E$.
\end{proposition}
\begin{proof}
First use exact positive-semidefinite and normalization tests to
validate the input.  Select independent columns $V_j$ of $E_j$ by
rational elimination.  For a nonzero effect its support is
\[
 P_j=V_j(V_j^*V_j)^{-1}V_j^*;
\]
for a zero effect use $P_j=0$.  Compress a rational Hermitian basis
by all $P_j$ and select a real-linearly independent subfamily
$W_1,\ldots,W_q$.  The positive Gram system
$\sum_b\tr(W_aW_b)c_b=\tr(W_aG)$ gives
$\Pi_{\mathcal W_E}G=\sum_bc_bW_b$.  This tests membership exactly.
Commutators test stationarity.  Formula~\eqref{eq:kerneldimension84}
then gives the dimension.  There are at most $kd^2$ candidate
matrices and $d^2$ independent matrices.  Clearing denominators
and applying fraction-free elimination gives polynomial bit
bounds by the determinant bounds used in
Proposition~\ref{prop:covarianceexact83}.  This is a decision about
supports and a specified tangent; no optimum over testers or
quantum-control synthesis is computed.
\end{proof}
